% ECONOMICS_PROBLEM: {"id": "G28-193", "title": "Liquidity regulation and solvency", "jel_code": "G28", "source_id": "OP-0321"}
% !TeX program = xelatex
\documentclass[11pt,a4paper]{article}
\usepackage[margin=23mm]{geometry}
\usepackage{fontspec}
\setmainfont{Latin Modern Roman}
\usepackage{amsmath,amssymb,mathtools}
\usepackage{xcolor,enumitem,booktabs,longtable,array,xurl,hyperref}
\definecolor{muted}{HTML}{53636C}
\hypersetup{colorlinks=true,urlcolor=blue,linkcolor=blue}
\setlength{\parindent}{0pt}
\setlength{\parskip}{5pt}
\newcommand{\R}{\mathbb R}
\newcommand{\E}{\mathbb E}
\newcommand{\N}{\mathbb N}
\newcommand{\ind}{\mathbf 1}
\newcommand{\Var}{\operatorname{Var}}
\newcommand{\Cov}{\operatorname{Cov}}
\newcommand{\supp}{\operatorname{supp}}
\DeclareMathOperator*{\argmin}{arg\,min}
\DeclareMathOperator*{\argmax}{arg\,max}
\newcommand{\entrymeta}[3]{{\sffamily\small\color{muted}\textbf{\texttt{#1}}\quad|\quad #2\par #3\par}}
\newcommand{\Problemref}[1]{\texttt{#1}}
\newcommand{\JELref}[2]{\texttt{#2} (JEL \texttt{#1})}
\begin{document}
\section*{Liquidity regulation and solvency}
\textbf{Problem ID:} \texttt{G28-193}\quad \textbf{JEL:} \texttt{G28}\par
% BEGIN ECONOMICS STATEMENT
\label{problem:OP-0321}
{\sffamily\small\color{muted}Original subject group: Banking, credit and financial stability\par}
\label{definitions:problem:30007603}
\textbf{Audit classification:} Published research agenda. The agenda addresses regulation interactions and resilience broadly. It does not supply the exact joint capital/liquidity objective or cross-partial. Those are editorial model choices, with local complementarity explicitly defined.
\subsection*{Definitions and precise research target}
Fix banks holding assets with specified durations, deposit liabilities and borrowing facilities. Let a joint rule \(a=(k,l)\) contain an equity requirement \(k\) and a liquid-asset requirement \(l\). A bank is solvent when its assets valued under the specified continuation allocation exceed liabilities; it is liquid when it can honor contemporaneous withdrawals without prohibited default. Interest-rate shocks change asset values, withdrawals affect financing needs and fire sales change market prices endogenously. Let \(W(a)\) be discounted social welfare in the resulting equilibrium, including intermediation services and crisis externalities. The task is to characterize the feasible joint optimum and its interaction term,
\[W_*=\sup_{a\in A}W(a),\qquad\chi=\frac{\partial^2W(k,l)}{\partial k\,\partial l},\]
where \(A\) is the legally and operationally feasible rule set and derivatives are interpreted only away from binding-regime boundaries. Identify whether capital and liquidity requirements are substitutes or complements in each state, and quantify changes when bank duration and deposit composition differ. Specify deposit-insurance and lender-of-last-resort arrangements because they affect both constraints. Resolve equilibrium feedbacks rather than imposing fixed fire-sale prices. Where the objective is non-smooth, report finite policy changes and an identified set of welfare-maximizing rules. A capital loss and a liquidity shortfall must not be treated as interchangeable outcomes.

\textbf{Additional formal definitions and scope (editorial).} Let liquidation value, continuation asset value and payment dates be distinct model objects. Solvency uses continuation value minus liabilities \(\geq0\) under a fixed valuation convention; liquidity uses cash plus permitted contemporaneous borrowing and liquidation proceeds sufficient for payments due. Regulatory \(k\) and \(l\) are ratios with explicitly specified denominators and asset eligibility; the deposit maturity and withdrawal technology are primitives. A positive cross-partial \(\chi\) means the marginal welfare benefit of raising \(k\) increases with \(l\), locally holding other rule components fixed. It is not a general statement that instruments are operational substitutes or complements across all objectives. If \(W\) is nonsmooth use the cross-difference \(W(k',l')-W(k',l)-W(k,l')+W(k,l)\) and report its steps. Welfare must include endogenous bank duration, funding and entry under each rule. For a nonempty feasible set and a finite optimal value, the design target is an \(\varepsilon\)-optimal feasible rule for each specified \(\varepsilon>0\): its cost is at most the infimum plus \(\varepsilon\), or its welfare is at least the supremum minus \(\varepsilon\). Report infeasibility or an unbounded value separately. An exact optimizer is requested only under additional attainment assumptions; it need not exist for the general rule class.
\subsection*{Variants and assumption changes}
\begin{enumerate}
\item \textbf{Fixed duration and insured deposits (editorial).} Estimate capital/liquidity welfare cross-differences with asset portfolios held technologically fixed.
\item \textbf{Endogenous duration and runs (editorial).} Banks choose asset duration and funding before shocks; compare joint rules under stated withdrawal and lender-of-last-resort assumptions.
\end{enumerate}
\subsection*{References and source audit}
\begin{enumerate}
\item Official January 9, 2026 web version; locator: Prudential Architecture / Regulatory framework and Resilience (broad framework only). Broad agenda; exact model editorial. URL: \url{https://www.bankofengland.co.uk/research/bank-of-england-agenda-for-research}.
\end{enumerate}
\begin{enumerate}\raggedright
\item\hypertarget{definitions:source:30007603:1}{} Bank of England. 2026. \emph{Bank of England Agenda for Research, 2025–2028}. Prudential Architecture / Regulatory framework and Resilience (broad framework only).
\url{https://www.bankofengland.co.uk/research/bank-of-england-agenda-for-research}.
\end{enumerate}

\subsection*{Common conventions: Source mathematical conventions}

These conventions apply to each entry unless explicitly replaced.

\textbf{Models and identification.} A primitive model \(m\) specifies all probability laws, feasible choices, information, equilibrium and measurement rules used by its target. The admissible class \(\mathcal M\) is fixed independently of outcomes. If \(P_O(m)\) is its observable probability law and \(\tau(m)\) the target, its identified set at observable law \(P\) is
\[
 \mathcal I_\tau(P)=\{\tau(m):m\in\mathcal M,\ P_O(m)=P\}.
\]
Point identification means that this set is a singleton. An empty set signals incompatibility of data and assumptions. Coordinate bounds need not describe the joint set. Bounds are sharp only if every retained target is attainable or arbitrarily approachable by compatible models. Finite bounds require explicit restrictions on unobserved quantities; unrestricted confounding, hidden wealth, extrapolation or arbitrary equilibrium selection can yield uninformative sets.

\textbf{Causal interventions.} A potential outcome \(Y_i(p)\) is the outcome under a complete policy regime \(p\), including timing, financing, information and market spillovers. The target population and units are fixed before treatment. With interference, use the complete assignment vector or a justified exposure mapping; individual no-interference notation alone cannot identify an economy-wide intervention. A difference of mean potential outcomes does not require the cross-world joint distribution. Conditioning on post-treatment migration, survival, enrollment or employment changes the estimand and needs separate justification.

\textbf{Statistical statements.} An estimator, confidence set, forecast or test specifies a sampling law, model class and loss/coverage criterion. Uniform \((1-\alpha)\)-coverage means \(\inf_{m\in\mathcal M}\Pr_m\{\tau(m)\in C_n\}\geq1-\alpha\), or an explicitly asymptotic version. The whole parameter space trivially has coverage, so informativeness requires a stated width, risk or power objective. A forecasting improvement is relative to a named benchmark, information set and evaluation population.

\textbf{Equilibrium and optimization.} An equilibrium comprises feasible plans, optimality, beliefs where needed, budget constraints and all market/resource clearing. The primitives or a stated benchmark define each correspondence \(\mathcal E(p;m)\). Nonemptiness is assumed or proved, never inferred just from compact policy sets. A selected-equilibrium target uses a declared selection rule; a robust target quantifies over all permitted equilibria. Use suprema or infima unless attainment is established. An \(\varepsilon\)-optimal policy is feasible and within \(\varepsilon\) of the finite optimum value. Report an empty feasible set separately.

\textbf{Welfare and accounting.} Normative utility, interpersonal normalization, social weights, numeraire, discounting and the use of public receipts are primitives. Behavioral utility need not coincide with the welfare criterion. Household welfare includes ownership income and rents once; adding profits again to compensating variation generally double-counts them. A monetary equivalent is defined by an explicit transfer and price convention, with monotonicity and range conditions. Average effects, mechanism contributions, transfers and welfare losses are different targets. Infinite sums require convergence and dynamic feasible plans require terminal or no-Ponzi conditions.

\textbf{Source versions.} A working-paper date, revision date and journal publication year are distinguished. A DOI for a journal version is not automatically the DOI of an earlier manuscript. Locators refer to the stated version and printed pages where specified. A metadata-only check does not verify a claim about a full-text conclusion. Every entry's source subsection narrows the provenance claim accordingly.
% END ECONOMICS STATEMENT
\end{document}
